\title{Parameter-Efficient Orthogonal Finetuning via Factorization}

\begin{abstract}
With the prevalence of large foundation models, the computational cost associated with training these models from the ground up has become unsustainable. Consequently, finding parameter-efficient strategies to adapt these sophisticated models for various downstream tasks is of increasing interest. In this work, we revisit Orthogonal Finetuning (OFT)—a principled strategy for task adaptation—highlighting its favorable generalization performance. However, OFT incurs considerable parameter overhead due to the dimensionality of orthogonal matrices. To mitigate this, we analyze OFT through the lens of information transmission and outline several criteria that facilitate higher parameter-efficiency. Drawing inspiration from the Cooley-Tukey fast Fourier transform, which streamlines information transfer, we propose utilizing butterfly structures for an efficient orthogonal parameterization. We integrate this approach within OFT, yielding a new parameter-efficient finetuning technique termed Orthogonal Butterfly~(BOFT). Notably, BOFT subsumes OFT as a particular case, thereby offering a broadened orthogonal finetuning framework. We thoroughly evaluate BOFT by applying it to finetune large-scale vision transformers, language models, and text-to-image diffusion models across a variety of visual and linguistic tasks.
\end{abstract}

\section{Introduction}

The impressive capabilities of recent models such as ChatGPT~\cite{brown2020language,chen2021evaluating} and Stable Diffusion~\cite{rombach2022high} exemplify the remarkable generalization strengths of large foundation models. Accompanying the notable leaps in model performance is a substantial escalation in parameter count (\emph{e.g.}, GPT-3 possesses about 175B parameters). As a result, the prospect of training such a model from scratch has become increasingly inaccessible for most practitioners. As broad advancement in the field now hinges on adapting existing foundation models rather than building them anew, efficient methodologies for transferring these models to new downstream tasks are indispensable. Predominant adaptation strategies include: \emph{model finetuning}~\cite{hulora2022,zaken2022bitfit,guo2020parameter,raffel2020exploring,qiu2023controlling,zhang2022adaptive,chavan2023one}, in which the model’s architecture is preserved and only certain parameters are adjusted; \emph{adapter tuning}~\cite{rebuffi2017learning,houlsby2019parameter,pfeiffer2020adapterfusion,lin2020exploring,he2021towards}, involving insertion of extra trainable components; and \emph{prompt tuning}~\cite{li2021prefix,lester2021power}, where learned prefix tokens are prepended to the model input. Among these, model finetuning stands out for its simplicity, effectiveness, and, notably, for avoiding any additional inference latency.

The main idea underpinning model finetuning is to maintain a close resemblance between the pretrained model and its finetuned counterpart, according to certain similarity metrics, thereby retaining the knowledge acquired during pretraining. For example, prevailing finetuning strategies typically utilize a small learning rate as a heuristic measure to keep modifications from pretraining to finetuning minimal. Considering that weight matrices are highly structured, a more systematic method focuses on preserving the relationships within these matrices—specifically, maintaining the angles between pairs of neurons. This perspective motivates the design of the Orthogonal Finetuning~(OFT)~\cite{qiu2023controlling} framework, in which all neurons in a given layer undergo transformation via the same orthogonal matrix, thereby guaranteeing that neuron pairwise angles remain invariant during finetuning. While OFT has shown significant improvements in both generalization and convergence when applied to finetuning text-to-image diffusion models~\cite{qiu2023controlling}, its reliance on large orthogonal matrices leads to a substantial number of trainable parameters. To make the parameter count more manageable, OFT adopts a block-diagonal format for these matrices. Yet, this increase in parameter efficiency introduces certain drawbacks: the fixed sparsity pattern of the orthogonal matrix constrains the flexibility of the transformation and enforces independence across blocks. This block-diagonal design, although beneficial for parameter reduction, may impose unwanted inductive biases (\emph{e.g.}, diminished expressivity and inability to approximate general linear maps), and crucially, the challenge of selecting an optimal sparsity configuration remains unresolved.

A central challenge in this context is constructing a dense orthogonal matrix without incurring excessive parameter counts. Although a standard $d$-dimensional dense orthogonal matrix typically involves $\mathcal{O}(d^2)$ parameters, our approach introduces an innovative method by synthesizing such matrices from multiple sparse factorized components. This methodology is motivated by the observation that computational efficiency can be sacrificed to reduce the parameter burden; specifically, expressing the orthogonal matrix as a sequence of sparse matrix products necessitates recurrent multiplication during every training step. Conceptually, this matrix decomposition can be reframed as an information transmission task: building a dense orthogonal matrix via sparse multiplicative factors is analogous to the process of propagating information through a grid-structured graph. Adopting this perspective enables the design of diverse sparse factorization strategies that strike a balance between model expressivity and parameter economy.

To optimize parameter utilization within this information transmission paradigm, we look to the butterfly network architecture employed in the Cooley-Tukey fast Fourier transform algorithm~\cite{cooley1965algorithm}, which efficiently enables information sharing among network nodes~\cite{klasing1994broadcasting}. The butterfly configuration intrinsic to the Cooley-Tukey procedure naturally yields a form of sparse matrix decomposition known as butterfly factorization. By leveraging butterfly factorization, a $d$-dimensional dense matrix can be realized as the product of $\mathcal{O}(\log d)$ sparse matrices, where each matrix contains only $\mathcal{O}(d)$ nonzero elements. When every constituent sparse matrix is additionally constrained to be orthogonal, we attain a substantially more efficient orthogonal parameterization—requiring just $\mathcal{O}(d\log d)$ parameters, as opposed to the conventional quadratic count. Utilizing this compact parameterization, we introduce a novel parameter-efficient finetuning approach, Orthogonal Butterfly~(BOFT), which generalizes and extends OFT as a specific instance within a broader orthogonal finetuning scheme. Both the block-diagonal and butterfly patterns share an underlying theme of sparsity, reducing the quantity of parameters that need to be learned while maintaining orthogonality. Notably, this sparsity-driven design philosophy stands in contrast to the low-rank parameterizations found in LoRA~\cite{hulora2022}; together, low-rank and sparse matrices represent fundamental structured matrix families~\cite{candes2011robust} widely adopted for parameter efficiency.

In contrast to the block-diagonal configuration employed by OFT—which strikes a balance between model expressiveness and regularization—BOFT adopts the butterfly structure to establish a more continuous transition between matrices representing the complete orthogonal group (capturing expressiveness) and the identity matrix (representing regularity). This design allows us to pinpoint a more compact subset of hypotheses within the orthogonal group that is suitable for downstream applications. Since the butterfly structure underlies numerous rapid linear transformations, including the discrete Fourier and discrete cosine transforms, we posit that our choice of structured parameterization not only enhances parameter efficiency but also naturally instills an inductive bias. This, in turn, aids generalization and helps mitigate overfitting. Our hypothesis is based on the observation that the butterfly structure is inherently capable of recovering a range of classical linear transformations, a feat not achievable by the block-diagonal structure used in OFT. The principal contributions of this work are summarized as follows:

\begin{itemize}[leftmargin=*,nosep]
    \item We introduce a fresh perspective on parameter efficiency within orthogonal finetuning, leveraging an information transmission framework that reinterprets the construction of parameter-efficient dense orthogonal matrices as an information flow problem over a grid-based graph.
    \item Drawing inspiration from the butterfly pattern prominent in the Cooley-Tukey algorithm, we develop Orthogonal Butterfly, a new method for orthogonal finetuning that achieves parameter efficiency within our information transmission framework.
    \item We offer several theoretical analyses elucidating how BOFT can sharply curtail the number of parameters that need to be trained, all while maintaining high expressivity and stable optimization. The underlying matrix factorization in BOFT further permits an insightful interpolation between weights (see Figure~\ref{fig:boft_interpolation}).
    \item To our knowledge, this work is the first to extend orthogonal finetuning to a range of domains beyond controllable text-to-image generation~\cite{qiu2023controlling}, demonstrating its strong promise as a general approach for model finetuning. Specifically, we showcase BOFT on a suite of adaptation tasks across computer vision and natural language processing, where it achieves notable performance gains over prevailing methods, thereby substantiating its parameter efficiency and robust generalization.
\end{itemize}

\section{Related Work}

\textbf{Parameter-efficient finetuning (PEFT)}. With the continuous growth in scale and capabilities of foundation models (\emph{e.g.}, \cite{brown2020language,radford2021learning,rombach2022high,kirillov2023segment}), there has been growing attention on methods to adapt these models to downstream applications without updating all parameters~\cite{treviso2023efficient,ding2023parameter,lialin2023scaling}. A wide range of PEFT approaches have been explored~\cite{houlsby2019parameter,aghajanyan2020intrinsic,hulora2022,edalati2022krona,wang2022adamix,gheini2021cross,zaken2022bitfit,guo2020parameter,sung2021training,ansell2022composable,lester2021power,li2021prefix,vu2022spot,he2021towards,mao2021unipelt,karimi2021compacter,liu2022few,sung2022lst,chen2022parameter,jia2022visual,chen2022adaptformer,zhang2022neural,jie2023fact,lian2022scaling,luo2023towards}, among which techniques that reparameterize network weights~\cite{aghajanyan2020intrinsic,hulora2022,edalati2022krona} are especially pertinent to our study. The LoRA framework~\cite{hulora2022}, for instance, modifies existing weight matrices by incorporating the product of two low-rank matrices, yielding notable outcomes on language tasks. Since LoRA fixes the rank of these low-rank matrices, several subsequent methods~\cite{zhang2022adaptive,valipour2022dylora,zhang2023increlora} have investigated adjusting ranks dynamically across layers to better distribute the available parameter budget. The ease of implementing this low-rank reparameterization has driven its widespread adoption~\cite{dettmers2023qlora,zi2023delta,chavan2023one}. Drawing on the relationship between hyperspherical energy and generalization~\cite{liu2018learning,liu2021orthogonal}, \cite{qiu2023controlling} introduce orthogonal finetuning (OFT) as a compelling alternative, wherein an orthogonal transformation is learned for neurons within the same layer, leading to improved generalization and more reliable training when compared to LoRA. Although OFT often delivers superior performance, it typically entails more trainable parameters relative to LoRA, highlighting the importance of developing more parameter-efficient forms of OFT. Furthermore, it remains uncertain whether OFT can be effectively extended to tasks beyond text-to-image diffusion model adaptation~\cite{qiu2023controlling}. BOFT addresses this by employing butterfly factorization to improve OFT's parameter efficiency. As a result, we can now showcase the benefits of orthogonal finetuning across a broader array of adaptation challenges.

\textbf{Butterfly structures}. The radix-2 Cooley-Tukey algorithm decomposes the $N$-point discrete Fourier transform into a sequence of smaller $\frac{N}{2}$-point discrete Fourier transforms through recursion, resulting in a butterfly structure—representable as the product of several sparse matrices, often collectively termed a butterfly matrix. These matrices have been exploited for various purposes, including representing orthogonal matrices to circumvent pivoting in Gaussian elimination and boost computational efficiency~\cite{parker1995random}, stabilizing recurrent neural network training~\cite{jing2017tunable}, and facilitating kernel approximations~\cite{munkhoeva2018quadrature}. Prior work has explored learning rapid linear transformations via butterfly parameterizations~\cite{dao2019learning,dao2019kaleidoscope}. Additionally, butterfly matrices have been employed for resource-efficient neural network training~\cite{chen2022pixelated,dao2022monarch}. Their utility further extends to speeding up matrix-vector products~\cite{michielssen1996multilevel,o2010algorithm}, enabling data-sparse matrix approximations~\cite{li2015butterfly}, and optimizing network transmissions~\cite{klasing1994broadcasting,li2003linear,rajasingh2016transmission}. Unlike these earlier applications, our investigation targets leveraging butterfly structures to render OFT more parameter-efficient.

\section{An Information Transmission View on Orthogonal Finetuning}

We begin by reviewing the fundamentals of OFT. When fine-tuning a pretrained weight matrix, OFT reformulates the updated weight matrix as a product between a trainable orthogonal matrix and the static pretrained matrix. Unlike LoRA—which augments weights via the addition of a low-rank matrix—OFT applies weight modifications through multiplicative orthogonal transformation. LoRA achieves parameter efficiency by adopting low-rank approximations, whereas classical OFT enforces a block-diagonal form on its orthogonal matrix~\cite{qiu2023controlling}; specifically, reducing the size of the diagonal blocks improves OFT's parameter efficiency. Figure~\ref{fig:comp_lora} provides a straightforward comparison between these approaches. The central rationale for using orthogonal transformations in fine-tuning is to maintain the angles between neurons~\cite{liu2018learning,liu2021orthogonal,qiu2023controlling}, which in turn helps preserve the semantic information learned during pretraining. In practice, OFT learns an orthogonal matrix $\bm{R}\in\mathbb{R}^{d\times d}$ for a given pretrained linear transformation $\bm{W}^0\in\mathbb{R}^{d\times n}$ and modifies the forward computation from $\bm{z}=(\bm{W}^0)^\top\bm{x}$ to $\bm{z}=(\bm{R}\bm{W}^0)^\top\bm{x}$, with $\bm{x}\in\mathbb{R}^{d}$ as input and $\bm{z}\in\mathbb{R}^{n}$ as output. OFT sets $\bm{R}$ to the identity matrix at initialization to guarantee zero deviation from the original weights. To maintain the orthogonality constraint on $\bm{R}$ during fine-tuning, we apply Cayley parameterization following~\cite{qiu2023controlling,liu2021orthogonal}: $\bm{R}=(\bm{I}+\bm{Q})(\bm{I}-\bm{Q})^{-1}$, where $\bm{Q}$ is skew-symmetric, i.e., $\bm{Q}=-\bm{Q}^\top$. The block-diagonal constraint for parameter savings is established by expressing $\bm{R}$ as $\text{diag}(\bm{R}_1,\bm{R}_2,\cdots,\bm{R}_r)$, where each $\bm{R}_i\in\mathcal{R}^{b\times b}$ is a compact orthogonal block, and $br=d$. However, while block-diagonality enhances parameter efficiency, it imposes the restriction that neuron dimensions (i.e., columns of $\bm{W}^0$) are partitioned into $r$ subsets, and each subset undergoes an independent orthogonal transformation. Although the block-diagonal approach has shown empirical utility, grouping neuron dimensions solely by index does not have a clear justification, suggesting that employing dense orthogonal matrices would be preferable. This prompts the following question: \emph{Is it possible to realize a dense orthogonal matrix structure while preserving the parameter efficiency?}

In response to this problem, we suggest constructing a dense orthogonal matrix by multiplying several sparse orthogonal matrices. To provide a cohesive and accessible way to analyze the sparsity patterns in orthogonal matrix factorizations, we reinterpret the construction of a dense orthogonal matrix within OFT as an information transmission scenario. Concretely, forming a dense matrix $\bm{R}\in\mathbb{R}^{d\times d}$ as a product of $m$ square matrices, $\thickmuskip=2mu \medmuskip=2mu \bm{R}=\bm{B}_m\bm{B}_{m-1}\cdots\bm{B}_1$, can be conceptualized as relaying information across a grid structure consisting of $d\times (m+1)$ nodes, as visualized in Figure~\ref{fig:info_trans}. The rationale behind this transmission framework lies in interpreting a $d$-dimensional dense square matrix as establishing complete connectivity from one set of $d$ nodes to another. In $\bm{R}$, if the entry $R_{ij}$ vanishes, information cannot travel from node $j$ to node $i$; a nonzero entry signifies that such transmission is possible. Therefore, representing the dense matrix $\bm{R}$ by the product $\bm{B}_m\bm{B}_{m-1}\cdots\bm{B}_1$ is analogous to propagating information step-by-step according to the network structures defined by each $\bm{B}_i$. Information first passes through the transformation given by $\bm{B}_1$, eventually reaching its destination after $\bm{B}_n$. 

To make this idea more concrete, Figure~\ref{fig:info_trans} illustrates the factorization $\bm{R}=\bm{B}_5\bm{B}_4\bm{B}_3\bm{B}_2\bm{B}_1$ alongside its respective sparsity patterns and the corresponding graph. Here, the depicted graph results from expanding the successive matrix multiplications. Each $\bm{B}_i$ can be viewed as defining the edges from the nodes at layer $i$ to those at layer $(i+1)$. In detail, the $(j_1,j_2)$ entry of $\bm{B}_i$ determines whether a directed connection exists from the $j_2$-th node in the $i$-th layer to the $j_1$-th node in the subsequent $(i+1)$-th layer (where a zero indicates no connection). For the overall product $\bm{B}_5\bm{B}_4\bm{B}_3\bm{B}_2\bm{B}_1$ to yield a dense matrix, it must enable every node in the first layer to relay information to every node in the sixth layer. 

If one examines only $\bm{R}=\bm{B}_4\bm{B}_3\bm{B}_2\bm{B}_1$, which refers to the connection between the first layer and the fifth, it is apparent that the first node cannot transmit to the third node, implying that the sparsity structure of $\bm{B}_4\bm{B}_3\bm{B}_2\bm{B}_1$ possesses a zero entry at $(3,1)$. This observation holds more generally: similar correspondences between the associated graphs and the resulting sparsity patterns arise when considering products such as $\bm{R}=\bm{B}_3\bm{B}_2\bm{B}_1$ or $\bm{R}=\bm{B}_2\bm{B}_1$. Broadly speaking, for any $\bm{R}\in\mathbb{R}^{d\times d}$ to be fully dense, the collection of matrices $\bm{B}_m,\cdots,\bm{B}_1$ must encode a set of directed links on a $d\times(m+1)$ node grid, where each directed edge only joins nodes in adjacent columns (i.e., levels), such that information is transmissible from every node in the initial layer to each node in the final $(m+1)$-st layer.

Adopting the information transmission perspective offers valuable insights into the sparse structure inherent in the factorization matrices of OFT. As illustrated in Figure~\ref{fig:blk_diag}, the original OFT constructs $\bm{R}$ with a block-diagonal arrangement. While this design efficiently reduces the count of parameters to be learned, it fails to realize a fully populated, or dense, matrix $\bm{R}$. The principal objective is to assemble a dense orthogonal matrix by combining $m$ sparse orthogonal matrices, simultaneously minimizing the count of significant, learnable parameters.

Through the information transmission lens, two main principles emerge: \emph{(i)} \textbf{dense connectivity}, meaning that every node at the first stage must connect via at least one route to every node in the final stage, and \emph{(ii)} \textbf{minimum free edges}, indicating that the aggregate number of edges should be minimized while upholding orthogonality. Orthogonality imposes specific restrictions on the permissible connections between adjacent layers. Notably, ensuring that each matrix $\bm{B}_i$ retains full rank—a prerequisite for orthogonality—necessitates $d$ edges to construct a one-to-one correspondence between all nodes at the $i$-th and $(i=1)$-th layers, so the number of inter-layer edges must not fall below $d$ (for example, see the case of $d=4$ in Figure~\ref{fig:info_trans}). As these $d$ required edges arise solely due to orthogonality, they are untrainable and, thus, should be omitted when tallying trainable parameters (as in the $d\times d$ orthogonal matrix with $d$ fixed $\pm1$ nonzeros). The property that orthogonal matrices require fewer parameters than their dense counterparts leads to added interdependence among connections: for example, if a $2\times2$ orthogonal matrix has a zero in the $(1,1)$ entry, a zero must also occur in the $(2,2)$ entry; thus, omitting one edge could enforce the absence of another. Consequently, orthogonality may either mandate or preclude certain edges in a given connection scheme. Examining the nonzero layout within the resultant orthogonal matrix, the information transmission viewpoint is especially advantageous for OFT, since the pattern and locations of trainable nonzeros in $\bm{R}$ are our focus, irrespective of their specific magnitudes.

A straightforward approach to establishing a dense connection between two layers requires $\mathcal{O}(d^2)$ connections, that is, utilizing a single dense orthogonal matrix, resulting in $d^2-d$ edges (for instance, with $d=4$, this amounts to 12 edges). As depicted in Figure~\ref{fig:info_trans}, a concrete example of matrix factorization is presented which uses just $10$ edges, fewer than a standard dense orthogonal matrix would. This approach allows us to analyze the parameter efficiency of OFT using a graphical model, enabling the construction of efficient factorizations via this representation. Our method is influenced by a notable connectivity scheme from the Cooley-Tukey algorithm, specifically the butterfly graph~\cite{cooley1965algorithm}, which achieves dense connections between $d$ input and $d$ output nodes with only $\mathcal{O}(d\log d)$ edges. For comparison, the structure shown in Figure~\ref{fig:info_trans} attains dense communication with $10$ edges, whereas the butterfly topology requires merely $8$ edges. In what follows, we describe how employing the butterfly structure furthers parameter efficiency.

\section{Orthogonal Parameterization by Butterfly Factorization}

Originally, the butterfly structure is integral to the Cooley-Tukey algorithm for fast Fourier transform. In the context of Fourier transform, modifications to frequencies affect the entire spatial representation, echoing the nature of our information transmission challenge—wherein each node at the initial layer can influence all nodes at the output layer. Additionally, the butterfly design is a favored topology in computer networking~\cite{solihin2015fundamentals,li2011linear} to enable efficient data dissemination. Given $k\geq 2$ as a power of $2$, we introduce the \emph{butterfly factor} $\bm{B}^F(k)$ by

\begin{equation}\label{eq:but_factor}
\begin{aligned}
    \bm{B}^F(k)=\begin{bmatrix}
    \text{diag}(\bm{d}_1) & \text{diag}(\bm{d}_2) \\
    \text{diag}(\bm{d}_3) & \text{diag}(\bm{d}_4)
    \end{bmatrix}\in\mathbb{R}^{k\times k},
\end{aligned}
\end{equation}

with each $\bm{d}_i\in\mathbb{R}^{\frac{k}{2}},\forall i$. Subsequently, for $d=2^N$, we recursively construct the $d$-dimensional \emph{butterfly matrix} $\bm{B}(d)\in\mathbb{R}^{d\times d}$ as follows:

\begin{equation}
\begin{aligned}
    \!\!\bm{B}(d)=\tilde{\bm{B}}(d,d)\!\cdot\!\begin{bmatrix}
    \bm{B}_1(\frac{d}{2})\!\!\!\!\! & \bm{0} \\
    \bm{0}\!\!\!\!\! &\bm{B}_2(\frac{d}{2})
    \end{bmatrix}= \tilde{\bm{B}}(d,d)\tilde{\bm{B}}(d,\frac{d}{2})\cdots\tilde{\bm{B}}(d,2),
\end{aligned}
\end{equation}

where $\bm{B}_1(\frac{d}{2})$ and $\bm{B}_2(\frac{d}{2})$ represent two butterfly matrices of dimension $\frac{d}{2}$. We introduce the \emph{butterfly component} as $\thickmuskip=2mu \medmuskip=2mu \tilde{\bm{B}}(d,k)=\text{diag}(\bm{B}^F_1(k),\cdots,\bm{B}^F_{d/k}(k))$, which constructs a block-diagonal $d\times d$ matrix where each block has size $k$, and each $\bm{B}^F_i(k)$, for all $i$, corresponds to a butterfly factor described in Equation~\ref{eq:but_factor}. This formulation enables us to employ butterfly matrices to represent orthogonal matrices. The crucial requirement is that all factors $\tilde{\bm{B}}(d,k)$, for every $k$, inside the overall butterfly matrix $\bm{B}(d)$ maintain orthogonality. Focusing on the particular case of the block-diagonal matrix $\tilde{\bm{B}}(d,2)$ (block size $2$), we can guarantee its orthogonality by parameterizing each $2\times 2$ block using the Cayley transform, or via $2$-dimensional rotation—similar to the techniques in \cite{liu2021orthogonal,qiu2023controlling}. Importantly, the sparsity structure (non-zero entries) of every butterfly component arises as a permutation of the sparsity pattern of $\tilde{\bm{B}}(d,2)$, which ensures that all butterfly components may be efficiently constructed as orthogonal matrices. As a result, this framework yields a scalable parameterization for orthogonal matrices, leveraging collections of $2\times 2$ orthogonal submatrices. Following \cite{chen2022pixelated}, we further extend the construction to define a block butterfly component, $\tilde{\bm{B}}^b(d,k)$, where every element of $\bm{d}_i$ (for each $i$) is replaced by a $b\times b$ matrix. To ensure that $\tilde{\bm{B}}^b(d,2)$ remains orthogonal, each block of dimension $2b\times 2b$ is parameterized as an orthogonal matrix. For larger blocks $\tilde{\bm{B}}^b(d,k)$ where $k>2$, the nonzero structure is obtained by permuting the blocks in the same pattern as $\tilde{\bm{B}}^b(d,2)$, and these can similarly be made orthogonal. Collectively, the BOFT forward computation is given by

\begin{equation}
    \begin{aligned}\nonumber
    \bm{z}=\big(\bm{R}(m,b)\cdot\bm{W}^0\big)^\top\bm{x},~~~~\text{s.t.}~\bigg{\{}\bm{R}(m,b)=\prod_{i=1}^m \tilde{\bm{B}}^b_{(d,i)}~~\&~~\underbrace{\big(\tilde{\bm{B}}^b_{(d,j)}\big)^\top\tilde{\bm{B}}^b_{(d,j)}=\tilde{\bm{B}}^b_{(d,j)}\big(\tilde{\bm{B}}^b_{(d,j)}\big)^\top\!=\bm{I}_d}_{\forall j\in [1, m]}\bigg{\}},
    \end{aligned}
\end{equation}

Here, for brevity, we define $\tilde{\bm{B}}^b(d,2^{m - i+1})$ as $\tilde{\bm{B}}^b_{(d,i)}$, and denote by $\bm{I}_d$ the $d$-dimensional identity matrix. The orthogonal matrix $\bm{R}(m,b)\in\mathbb{R}^{d\times d}$ is constructed through the multiplication of several orthogonal butterfly modules. To streamline notation, we represent BOFT configured with $\bm{R}(m,\frac{b}{2})$ as BOFT($m$,$b$), for all $b\geq 2$. Setting $\thickmuskip=2mu \medmuskip=2mu m=1$, BOFT($1$,$b$) simplifies to the block-diagonal OFT~\cite{qiu2023controlling} with blocks of size $b$, while BOFT($1$,$d$) becomes the conventional OFT~\cite{qiu2023controlling} featuring an unconstrained full orthogonal matrix. Moreover, BOFT($\log \frac{2d}{b}$,$b$) takes the block butterfly matrix $\bm{B}^{\frac{b}{2}}(d)$ for $\bm{R}$, resulting in a dense orthogonal matrix $\bm{R}$. In the general case, the parameter count for BOFT($m$,$b$) amounts to $\frac{1}{2}(b-1)dm$ effective trainable entries when adapting a linear transformation of size $d\times n$. If one adopts the standard butterfly construction with $m=\log d, b=2$, BOFT requires only $\mathcal{O}(d\log d)$ parameters. This contrasts sharply with the classic OFT employing a fully dense orthogonal matrix, which demands $\mathcal{O}(d^2)$ parameters, and the block-diagonal OFT with $r$ blocks, which requires $\mathcal{O}(bd)$. Consequently, producing a dense orthogonal matrix using the original OFT necessitates selecting $b=d$, while BOFT allows for any choice of $b$ to achieve density.

\textbf{BOFT Identity Initialization}. Common finetuning approaches typically commence with the weights from the pretrained model, aiming to avoid substantial divergence from the original parameters. Analogous to LoRA's use of zero initialization for low-rank modifications, BOFT begins by setting all butterfly factors as identity matrices (i.e., initializing the skew-symmetric part to zero within the Cayley transform).

\textbf{BOFT Multiplicative Dropout}. LoRA~\cite{hulora2022} enhances generalization by integrating Dropout into its low-rank adaptation modules. While standard Dropout~\cite{srivastava2014dropout} operates seamlessly for LoRA's additive framework, it is not compatible with BOFT's multiplicative parameterization. To overcome this limitation, we devise a multiplicative Dropout tailored to BOFT. Since $\bm{R}(m,b)$ consists of $m$ orthogonal butterfly blocks, which can be rearranged into $2b\times 2b$ block-diagonal orthogonal matrices, the multiplicative Dropout method stochastically selects $p_1$ percent of butterfly components and $p_2$ percent of diagonal sub-blocks per component, replacing the chosen segments with identity matrices.

\textbf{Expressivity of BOFT}. While the butterfly structure paired with permutations can exactly implement a range of well-known fast linear transforms~\cite{dao2019learning,dao2019kaleidoscope} (\emph{e.g.}, fast Fourier transform, Hadamard transform), the capacity of our orthogonal butterfly matrix to approximate arbitrary orthogonal matrices remains uncertain. To investigate this, we simulate the approximation of a randomly sampled dense orthogonal matrix~\cite{anderson1987generation} of dimension $512\times 512$. Results are presented in Figure~\ref{fig:para_eff}, averaged over 10 distinct random initializations. Here, the y-axis corresponds to the measured approximation error, while the x-axis reflects the number of effective trainable parameters. Each colored curve represents BOFT for a fixed block size; the curve’s leftmost marker corresponds to the error for BOFT($1$,$b$) (\emph{i.e.}, the conventional block-diagonal OFT using block size $b$). Overall, BOFT achieves superior parameter efficiency in comparison to OFT. For instance, BOFT($9$,$2$) not only attains higher expressiveness than BOFT($1$,$16$), but does so with significantly fewer parameters. Increasing $m$ while decreasing $b$ generally leads to greater parameter efficiency; as an example, BOFT($6$,$4$) matches the approximation performance of BOFT($2$,$16$) with considerably fewer parameters. The butterfly structure, by virtue of its richer configuration space relative to block-diagonal forms, renders BOFT strictly more expressive than OFT for a given block size.

\begin{theorem}[Expressivity of BOFT]\label{thm:exp_boft}
    BOFT possesses greater expressiveness compared to OFT with the same block size. Any orthogonal matrix of dimension $d$ can be approximated using products of butterfly matrices as $\bm{B}_{d-1,1}(d)\bm{B}^\top_{d-1,2}(d)\cdots\bm{B}_{1,1}(d)\bm{B}^\top_{1,2}(d)$, where every $\bm{B}_{i,j}(d)$ is a butterfly matrix.
\end{theorem}

Theorem~\ref{thm:exp_boft} naturally inspires an extension of BOFT, where the overall orthogonal transformation takes the form $\thickmuskip=2mu \medmuskip=2mu \bm{R}^G(m_1,b_1,m_2,b_2,l)=\bm{R}_{l,1}(m_1,b_1)\bm{R}_{l,2}^T(m_2,b_2)\cdots\bm{R}_{1,1}(m_1,b_1)\bm{R}_{1,2}^T(m_2,b_2)$, with $\bm{R}_{i,j}^T(m,b)$ representing the constituent orthogonal matrices in BOFT. Setting $m_1=m_2=\log d$, $b_1=b_2=2$, and $l=d-1$, this construction $\bm{R}^G(m_1,b_1,m_2,b_2,l)$ can realize any element from the orthogonal group—this hierarchical organization is also referred to as the kaleidoscope hierarchy~\cite{dao2019kaleidoscope}. Nevertheless, increased expressivity does not inherently ensure enhanced performance during finetuning; for example, full finetuning, though universally expressive, often leads to suboptimal results. Striking an optimal balance between expressivity and regularity is essential for model finetuning generalization. BOFT expands the space of finetunable parameters while retaining structured inductive biases, facilitating an improved trade-off between expressivity and regularization.

\textbf{Spectral properties}. In terms of maintaining spectral characteristics, orthogonal finetuning outperforms LoRA as it exactly conserves the spectral norm of the pretrained weight matrix $\bm{W}^0$. This can be understood via singular value decomposition: $\bm{W}^0 = \bm{U}\bm{\Sigma}\bm{V}^\top$, with $\bm{U}, \bm{V}$ being orthogonal matrices and $\bm{\Sigma}$ a diagonal matrix of singular values. Both OFT and BOFT operate by left-multiplying an orthogonal matrix $\bm{R}$, producing updated weights $\bm{R}\bm{U}\bm{\Sigma}\bm{V}^\top$, an operation that leaves the principal singular value, and thus the spectral norm of $\bm{W}^0$, unchanged. Maintaining the spectral norm has been empirically linked to improved training robustness and better generalization~\cite{miyato2018spectral,yoshida2017spectral}. Further mathematical aspects are detailed in Appendix~\ref{app:sn_property}.

\textbf{Orthogonal finetuning as learning bilinear similarity}. The BOFT method can be reformulated as the optimization of a bilinear similarity function $\bm{w}^0_i\bm{R}\bm{x}$, where $\bm{w}^0_i$ denotes the $i$-th column of the pretrained matrix $\bm{W}^0$. In this view, BOFT seeks to optimize the matrix $\bm{R}$—subject to the orthogonality constraint—as a bilinear similarity matrix. This ties closely to established concepts in distance metric learning~\cite{xing2002distance} and the theory of bilinear forms~\cite{roman2005advanced}.

\textbf{Inductive bias and generalization in BOFT}. The parameter $\bm{R}(m,b)$ within the BOFT approach typically embodies a structured element of the orthogonal group, thereby restricting the space of candidate models. As a result, BOFT instills a particular inductive bias, which we posit is advantageous for generalization. This claim rests on the observation that the structured patterns generated by butterfly factorization closely mirror those of several fundamental linear transforms~\cite{dao2019kaleidoscope}, such as the discrete Fourier, sine/cosine, and Hadamard transforms. Additionally, the use of sparse matrix factorization in BOFT may further contribute beneficial, implicit inductive biases~\cite{gunasekar2017implicit,li2018algorithmic,liu2019neural}.

\textbf{Comparison to butterfly-based sparse training}. Numerous studies~\cite{dao2019learning,dao2019kaleidoscope,chen2022pixelated,dao2022monarch} have explored sparse training methods grounded in the butterfly parameterization. These works predominantly address the direct reparameterization of weight matrices via butterfly structures, training neural networks ab initio. Notably, \cite{dao2022monarch} examines a finetuning technique where pretrained weights are first projected onto a modified form of butterfly matrices before optimizing the projected elements for downstream applications. In contrast, BOFT introduces a markedly different approach to finetuning by applying layer-wise shared weight transformations.

\input{rebuttal-results}

\section{Concluding Remarks and Limitations}

In this work, we introduce Orthogonal Butterfly, a broadly applicable and parameter-efficient finetuning approach for foundation models leveraging the butterfly structure. The principal contribution is the realization that representing a dense orthogonal matrix as a product of several sparse orthogonal matrices significantly enhances parameter efficiency. To facilitate the discovery of suitable matrix factorizations, we propose a graphical information transmission framework. Within this framework, the butterfly construction is identified as a promising mechanism for achieving efficient sparse orthogonal decompositions. We empirically validate the effectiveness of BOFT for finetuning various large-scale models, including large language models, large vision models, and text-to-image generative models, demonstrating its advantages as a universal finetuning strategy.

However, BOFT still has certain limitations. As the method constructs the final orthogonal matrix by sequentially composing multiple orthogonal matrices, its training phase incurs somewhat greater computational overhead compared to OFT. Addressing ways to reduce this training cost in BOFT is a direction for future investigation. Encouragingly, the orthogonal matrices obtained via BOFT can be merged into the original pretrained models after finetuning, incurring no extra inference overhead. Additionally, it is not yet clear whether the butterfly network structure achieves optimal efficiency for information transmission. Our graphical information transmission framework enables cross-pollination with fields such as computer networking, where the study of network topology for information transfer is extensive. We anticipate that future research might identify even more efficient network architectures for assembling dense orthogonal matrices.

\end{document}